% Gerado por src/python/thesis/build_tab_magnitude_confirmatoria.py -- nao editar a mao.
% Regenerar com: make thesis-tables
% Fonte: data/processed/todas_metricas_consolidado_with_modern.csv
\begin{table}[!htb]
    \centering
    \small
    \setlength{\tabcolsep}{5pt}
    \caption{Sensibilidade exploratória da diferença relativa das medianas de IGD.}
    \label{tab:magnitude_sensibilidade}
    \begin{tabular*}{\textwidth}{@{\extracolsep{\fill}}llrr@{}}
        \toprule
        \textbf{Instância} & \textbf{$M$} & \textbf{$\Delta$ IGD (\%)} & \textbf{IC bootstrap 95\%} \\
        \midrule
        WFG1 & 2 & $+24{,}07$ & [$-4{,}55$, $+37{,}34$] \\
        WFG6 & 2 & $-8{,}87$ & [$-25{,}55$, $+3{,}81$] \\
        DTLZ4 & 3 & $+0{,}92$ & [$-884{,}49$, $+2{,}07$] \\
        MaF5 & 3 & $-0{,}21$ & [$-505{,}66$, $+1{,}66$] \\
        \bottomrule
    \end{tabular*}
    \par\smallskip
    \parbox{\textwidth}{\footnotesize
    Coorte: IVF/SPEA2 nas execuções 3001--3060 contra SPEA2 nas execuções 1--60, 60 execuções independentes por algoritmo e instância. $\Delta$ é a diferença entre medianas em porcentagem da mediana do SPEA2; valores positivos favorecem IVF/SPEA2. O IC percentil de 95\% resulta de 10.000 reamostragens independentes com reposição dentro de cada algoritmo/instância (semente 20260927). Ele não é ajuste de Holm, não demonstra equivalência e não é intervalo para uma proporção de problemas; a tabela não sustenta inferência global.
    }
\end{table}
